\documentclass[11pt,a4paper]{scrartcl}
\usepackage{iutparakou}
\begin{document}

\fichetitre{Jour 4 — Calculabilité}{Machine de Turing \& machine universelle}
{Modules : \texttt{turing.py}, \texttt{utm.py} \quad$\bullet$\quad App : \texttt{apps/mtu} \quad$\bullet$\quad Barème : 4 pts (+1)}

\section*{Contexte}
On lève la dernière limite : une mémoire \textbf{libre}, lisible \emph{et}
ré-inscriptible n'importe où — la machine de Turing. Puis une idée vertigineuse :
\textbf{une seule} machine peut simuler \textbf{toutes} les autres si on lui donne leur
description. C'est la machine \textbf{universelle}, ancêtre du processeur et de l'interpréteur.

\section*{Notions nouvelles du jour}

\subsection*{1. La machine de Turing}
$M=(Q,\Sigma,\Gamma,\delta,q_0,\text{\textvisiblespace},F)$ avec
$\delta:Q\times\Gamma\to Q\times\Gamma\times\{L,R,S\}$. Une \textbf{tête} lit/écrit une
case et se déplace : « dans l'état $q$, je lis $a$ $\to$ j'écris $b$, je bouge de $d$, je
passe en $q'$ ». Différence avec la pile (jour~2) : ici on peut \textbf{relire et
réécrire} le passé.

\begin{center}
\begin{tikzpicture}[start chain=going right, node distance=0mm]
  \foreach \c/\i in {\_/-1,1/0,1/1,+/2,1/3,\_/4,\_/5} {\node[tapecell, on chain] (c\i) {\c};}
  \node[above=5mm of c0, iutorange!85!black] (tete) {\small\bfseries tête, état \texttt{q0}};
  \draw[tapehead] (tete) -- (c0);
  \node[right=6mm of c5, font=\small\itshape, text width=30mm] {ruban bi-infini ; blanc \texttt{\_} ; entrée \texttt{11+1}.};
\end{tikzpicture}
\end{center}

\subsection*{2. Une opération \emph{est} une machine : trace de \texttt{ADD} sur \texttt{11+1}}
Table : \texttt{(q0,1)$\to$(q0,1,R)} ; \texttt{(q0,+)$\to$(q1,1,R)} ;
\texttt{(q1,1)$\to$(q1,1,R)} ; \texttt{(q1,\_)$\to$(q2,\_,L)} ; \texttt{(q2,1)$\to$(qf,\_,S)}.
La case sous la tête est notée \texttt{[\,]}.

\begin{exemple}[Trace pas-à-pas — \texttt{2+1 = 3}]
\begin{center}\small\ttfamily
\begin{tabular}{@{}clll@{}}
\toprule
\rmfamily pas & \rmfamily état & \rmfamily ruban & \rmfamily action\\\midrule
0 & q0 & [1]1+1 & lit 1 $\to$ R\\
1 & q0 & 1[1]+1 & lit 1 $\to$ R\\
2 & q0 & 11[+]1 & \rmfamily le \texttt{+} devient \texttt{1}, R, passe q1\\
3 & q1 & 111[1] & lit 1 $\to$ R\\
4 & q1 & 1111[\_] & \rmfamily blanc : demi-tour, passe q2\\
5 & q2 & 111[1]\_ & \rmfamily efface le dernier 1, passe qf\\
6 & qf & 111[\_] & \rmfamily \bfseries accepté ; ruban = 111 = 3\\
\bottomrule
\end{tabular}
\end{center}
Idée : \texttt{+}$\to$\texttt{1} \emph{fusionne} les deux blocs (on a $n{+}1{+}m$ uns),
puis on retire \textbf{un} \texttt{1}. Coût $O(n{+}m)$.
\end{exemple}

\subsection*{3. Lire une table par \emph{phases} : la soustraction}
\texttt{SUB} apparie un \texttt{1} de $m$ (marqué \texttt{X}) avec un \texttt{1} de $n$
(barré \texttt{X}, de droite à gauche), jusqu'à épuiser $m$, puis nettoie. Chaque état
est une \textbf{phase} :

\begin{center}
\begin{tikzpicture}[node distance=20mm and 22mm, on grid]
  \node[etat] (q0) {q0};
  \node[etat, right=of q0] (q1) {q1};
  \node[etat, right=of q1] (q2) {q2};
  \node[etat, below=of q1] (q3) {q3};
  \node[etat, right=22mm of q2] (q4) {q4};
  \node[etatfin, right=of q4, accepting] (qf) {qf};
  \path[fleche]
    (q0) edge node[above]{\scriptsize \texttt{-} : R} (q1)
    (q1) edge node[above]{\scriptsize \texttt{1$\to$X}, L} (q2)
    (q2) edge[bend left=15] node[below,pos=.4]{\scriptsize \texttt{-} : L} (q3)
    (q3) edge[bend left=15] node[above,pos=.6]{\scriptsize \texttt{1$\to$X}, R} (q0)
    (q1) edge[bend left=20] node[above]{\scriptsize \texttt{\_} : $m$ épuisé} (q4)
    (q4) edge node[above]{\scriptsize nettoie X,\texttt{-}} (qf);
\end{tikzpicture}
\end{center}
\noindent \textbf{Trois points délicats} (à expliquer au rapport) : \texttt{(q0,X)$\to$R}
et \texttt{(q1,X)$\to$R} — la table doit gérer ses \emph{propres} marques ; le barrage
\textbf{droite$\to$gauche} donne une sortie contiguë ; la fin de $m$ est détectée par un
\textbf{blanc} (le ruban compte, pas un compteur).

\subsection*{4. Codage $\langle M\rangle$ et machine universelle}
Une MT n'est qu'une \textbf{table finie} : on l'écrit comme une \textbf{chaîne}
$\langle M\rangle$ (\textbf{injective}, \textbf{décodable}). Un \emph{programme} devient
une \emph{donnée}. La machine \textbf{universelle} $U$ prend $\langle M\rangle\,\#\#\,w$
et \textbf{simule} $M$ (principe du \emph{programme enregistré}). Un \textbf{cycle} :

\begin{center}
\begin{tikzpicture}[noeudc/.style={draw=iutblue, very thick, rounded corners=3pt,
    fill=iutblue!6, minimum height=8mm, inner xsep=5pt, font=\small\bfseries, text=iutblue},
    fl/.style={-{Stealth[length=2mm]}, very thick, iutorange!85!black}]
  \node[noeudc] (lire) {1. Lire le symbole simulé};
  \node[noeudc, right=12mm of lire] (chercher) {2. Chercher la transition dans $\langle M\rangle$};
  \node[noeudc, below=8mm of chercher] (ecrire) {3. Écrire le symbole};
  \node[noeudc, left=12mm of ecrire] (bouger) {4. Déplacer la tête};
  \node[noeudc, below=8mm of lire] (maj) {5. Mettre à jour l'état};
  \draw[fl] (lire) -- (chercher); \draw[fl] (chercher) -- (ecrire);
  \draw[fl] (ecrire) -- (bouger); \draw[fl] (bouger) -- (maj); \draw[fl] (maj) -- (lire);
\end{tikzpicture}
\end{center}

\begin{notion}[Church--Turing \& composition]
« Toute fonction effectivement calculable est calculable par une MT » : pas un théorème,
mais une thèse étayée par l'\textbf{équivalence} de modèles indépendants ($\lambda$-calcul,
fonctions récursives, RAM). On obtient $\times$ et $\div$ par \textbf{composition} de MT
(mult.\ = additions répétées) — ce que fait $U$ en orchestrant des sous-programmes
(\texttt{apps/mtu}). La calculatrice est \textbf{exécutée par} $U$ (\texttt{addition\_via\_utm}).
\end{notion}

\subsection*{5. La limite : indécidabilité}
Revers de l'universalité : \textbf{on perd la décidabilité}. Le problème de l'\textbf{arrêt}
(\textsc{halt}) est indécidable (Turing 1936) ; par réduction, « $U$ s'arrête sur
$\langle M\rangle\#\#w$ » l'est aussi. \emph{Universel $\neq$ omniscient} : $U$ peut tout
\textbf{simuler}, mais aucune machine ne \textbf{prédit} en général l'arrêt d'une autre.

\section*{A. Théorie \normalfont(à rédiger) — 2 pts}
\begin{description}[leftmargin=1.6em,style=nextline]
  \item[Q4.1 (0,5)] Schéma d'encodage $\langle M\rangle$ ; justifier injectif + décodable.
  \item[Q4.2 (0,5)] Décrire $U$ et son cycle de simulation (invariant + un tour).
  \item[Q4.3 (0,5)] Surcoût (\textbf{sourcé}) : multi-ruban $O(t\cdot|\langle M\rangle|)$ ;
    multi$\to$mono $\le$ quadratique ; Hennie--Stearns (logarithmique).
  \item[Q4.4 (0,5)] Indécidabilité par réduction depuis \textsc{halt}.
\end{description}

\section*{B. Pratique — 2 pts}
\begin{description}[leftmargin=1.6em,style=nextline]
  \item[E4.1 (0,5)] \texttt{TuringMachine.run} (ruban, arrêt, \texttt{trace}).
  \item[E4.2 (0,5)] \texttt{encode}/\texttt{decode} + \texttt{UniversalTM.run}.
  \item[E4.3 (0,75)] Tables \texttt{ADD}/\texttt{SUB} + \texttt{Calculatrice} ($\times,\div$ par composition).
  \item[E4.4 (0,25)] \texttt{UniversalInterpreter.run} (encode puis délègue à $U$).
\end{description}

\begin{gate}
\texttt{apps/mtu} ne réécrit aucune boucle de MT : il décrit des \textbf{tables} et les
fait tourner via \texttt{formlang.turing}/\texttt{utm}. \emph{Honnêteté} : surcoût
\textbf{référencé}, pas de constante « de mémoire ».
\end{gate}

\subsection*{Références}
Turing 1936. \textbullet\ Sipser (MTU, indécidabilité). \textbullet\ Hennie \& Stearns
1966 (\emph{J.~ACM} 13(4)). \textbullet\ Arora \& Barak 2009 (surcoûts).

\end{document}
